\chapter{Entorno de software y primera prueba}
\label{chap:software}

\section{Preparar el entorno}

BrainChip distribuye las herramientas de Akida mediante MetaTF. En la versión utilizada
para este manual, la instalación actual emplea \code{metatf==2.19.3} y admite
Python 3.10--3.12 \cite{brainchip_install_2193}.

El repositorio también incluye una configuración de referencia basada en Python 3.11
y Akida 2.19.1, la misma versión utilizada durante el TFG. Esta segunda opción permite
repetir la prueba básica con un entorno conocido.

Abre una terminal en \projectroot, la carpeta extraída del ZIP.

En Linux:

\begin{lstlisting}[language=bash]
python3.11 -m venv .venv
source .venv/bin/activate
python -m pip install --upgrade pip
\end{lstlisting}

En Windows, PowerShell:

\begin{lstlisting}
py -3.11 -m venv .venv
.\.venv\Scripts\Activate.ps1
python -m pip install --upgrade pip
\end{lstlisting}

El entorno virtual solo se crea una vez. En una terminal nueva basta con volver a
activarlo antes de continuar.

\section{Instalar Akida}

Elige una de las dos configuraciones siguientes.

Para trabajar con la versión actual documentada por BrainChip:

\begin{lstlisting}
python -m pip install "metatf==2.19.3"
\end{lstlisting}

Para repetir la configuración de referencia incluida con el repositorio:

\begin{lstlisting}
python -m pip install -r requirements/reproducible.txt
\end{lstlisting}

No mezcles ambas instalaciones en el mismo entorno virtual.

A continuación, instala las herramientas incluidas con este repositorio:

\begin{lstlisting}
python -m pip install -e . --no-deps
python -m pip check
\end{lstlisting}

Comprueba que Akida se importa correctamente:

\begin{lstlisting}
python -c "import akida; print(akida.__version__); print(akida.devices())"
\end{lstlisting}

Si no hay una tarjeta conectada, es normal que \code{akida.devices()} devuelva
\code{[]}. Si el equipo tiene una AKD1000 conectada y el controlador está funcionando,
debe aparecer al menos un dispositivo.

\section{Primera prueba en software}

Antes de utilizar la tarjeta, ejecuta el modelo de prueba incluido con el repositorio:

\begin{lstlisting}
python examples/00_simulator_smoke.py
\end{lstlisting}

El modelo recibe cuatro valores de entrada. Para

\begin{lstlisting}
[3, 2, 1, 4]
\end{lstlisting}

la salida esperada es

\begin{lstlisting}
[5, 5]
\end{lstlisting}

El script ejecuta primero el modelo en software y comprueba después que puede mapearse
sobre un AKD1000 virtual. Este último paso verifica la compatibilidad del modelo con la
arquitectura de la tarjeta, pero no utiliza hardware físico.

También puedes ejecutar la comprobación incluida en la herramienta del repositorio:

\begin{lstlisting}
akd1000-bringup verify --id smoke --model models/akida_v1_smoke_test.fbz
\end{lstlisting}

Si ambas pruebas terminan correctamente, el entorno de Python y el modelo de prueba
están preparados.

\section{Comprobar el equipo}

Antes de usar la tarjeta física, puedes ejecutar el diagnóstico incluido con el
repositorio:

\begin{lstlisting}
akd1000-bringup doctor
\end{lstlisting}

Para guardar el resultado:

\begin{lstlisting}
mkdir -p evidence
akd1000-bringup doctor --json > evidence/preflight.json
\end{lstlisting}

La herramienta comprueba, entre otros datos, la versión de Akida, el kernel, la conexión
PCIe y los dispositivos detectados. Su salida ayuda a localizar el problema si la tarjeta no aparece.

En un equipo sin AKD1000 es normal que el diagnóstico indique que no se ha encontrado
ningún dispositivo.

\section{Primera ejecución en la AKD1000}

Esta prueba requiere que la tarjeta haya sido detectada por PCIe y que el controlador
del \cref{chap:host-driver} esté instalado.

Comprueba primero que Akida detecta el dispositivo:

\begin{lstlisting}
akida devices
\end{lstlisting}

Después ejecuta:

\begin{lstlisting}
python examples/01_hardware_smoke.py
\end{lstlisting}

El script carga el mismo modelo utilizado en la prueba de software, lo mapea sobre la
AKD1000 y ejecuta la entrada conocida. La salida debe volver a ser:

\begin{lstlisting}
[5, 5]
\end{lstlisting}

Un resultado \code{PASS} confirma que el modelo se ha cargado, que la tarjeta es
accesible desde Akida y que la inferencia se ha ejecutado correctamente en el hardware.

\paragraph{Si la prueba física falla.}
No continúes todavía con un modelo propio. Conserva el mensaje de error y utiliza el
procedimiento de diagnóstico del \cref{chap:diagnostico} para localizar el problema.
